\documentclass[11pt]{article}
\usepackage{reusv}
\usepackage{listings}
\usepackage{enumitem}

\lstset{
  basicstyle=\ttfamily\small,
  frame=single,
  breaklines=true,
  language=Python,
  showstringspaces=false,
  commentstyle=\color{gray},
  keywordstyle=\color{blue},
  stringstyle=\color{red}
}
\setborderthickness{0.5pt}
\setbordermargin{1cm}
\setbordercolor{black}

\slimtitle{D1. 가우시안 프로세스 분류기}

\begin{document}
\drawborder
\thispagestyle{firstpage}

\lightertitle{D1. 가우시안 프로세스 분류기}
\def\lighterdate{February 12, 2026}
\lighterauthor{Suwon Lee}
\lighteraddress{}{Future Mobility Control Lab, Department of Future Mobility, Kookmin University, Seoul, Republic of Korea}
\lighteremail{suwon@kookmin.ac.kr}

\begin{abstract}
본 프로젝트는 LEO-to-LEO 연속 추력 궤도전이의 최적 추력 프로파일을 3가지 유형(unimodal, bimodal, multimodal)으로 분류한다.
분류의 입력은 3차원 정규화 파라미터 공간 $\mathbf{x} = [\tilde{T},\, \Delta a,\, \Delta i]^\top \in [0,1]^3$이며,
각 점에서의 클래스 레이블은 direct collocation으로 산출한 최적 추력 프로파일의 피크 수로 결정된다.
이 분류 문제를 가우시안 프로세스 분류기(GPC)로 모델링한다.
GPC는 결정 경계에서의 예측 불확실성을 확률적으로 정량화하므로,
능동 학습에서 다음 샘플링 위치를 결정하는 acquisition function의 기반이 된다.
본 모듈(\texttt{gp\_classifier.py})은 scikit-learn의 \texttt{GaussianProcessClassifier}를 래핑하여,
ARD RBF 커널과 Laplace 근사 기반의 3-class GPC를 제공한다.
\end{abstract}


\section{목적}

본 프로젝트는 LEO-to-LEO 연속 추력 궤도전이의 최적 추력 프로파일을 3가지 유형으로 분류한다.

\begin{table}[h]
\centering
\begin{tabular}{ccl}
\toprule
클래스 $c$ & 레이블 & 물리적 의미 \\
\midrule
0 & Unimodal & 단일 추력 arc (피크 1개) \\
1 & Bimodal & 이중 추력 arc (피크 2개) \\
2 & Multimodal & 다중 추력 arc (피크 3개 이상) \\
\bottomrule
\end{tabular}
\end{table}

분류의 입력은 3차원 정규화 파라미터 공간 $\mathbf{x} = [\tilde{T},\, \Delta a,\, \Delta i]^\top \in [0,1]^3$이다.
여기서 $\tilde{T}$는 정규화 추력, $\Delta a$는 반장축 변화량, $\Delta i$는 궤도경사 변화량이다.
각 점에서의 클래스 레이블은 direct collocation으로 산출한 최적 추력 프로파일의 피크 수로 결정된다.

이 분류 문제를 가우시안 프로세스 분류기(Gaussian Process Classifier, GPC)로 모델링한다~\cite{Rasmussen2006}.
GPC는 결정 경계에서의 예측 불확실성을 확률적으로 정량화하므로,
능동 학습(active learning)~\cite{Settles2009}에서 다음 샘플링 위치를 결정하는 acquisition function의 기반이 된다.
본 모듈(\texttt{gp\_classifier.py})은 scikit-learn~\cite{Pedregosa2011}의 \texttt{GaussianProcessClassifier}를 래핑하여,
ARD(Automatic Relevance Determination) RBF 커널과 Laplace 근사 기반의 3-class GPC를 제공한다.


\section{수학적 배경}

\subsection{가우시안 프로세스 회귀의 기초}

\textbf{가우시안 프로세스(Gaussian Process, GP)}는 함수에 대한 확률 분포이다~\cite{Rasmussen2006}.
GP는 평균 함수 $m(\mathbf{x})$와 공분산 함수(커널) $k(\mathbf{x}, \mathbf{x}')$로 완전히 정의된다:
\[
f \sim \mathcal{GP}\bigl(m(\mathbf{x}),\; k(\mathbf{x}, \mathbf{x}')\bigr)
\]

이것은 임의의 유한 입력 집합 $\{\mathbf{x}_1, \dots, \mathbf{x}_N\}$에 대해,
함수값 벡터 $\mathbf{f} = [f(\mathbf{x}_1), \dots, f(\mathbf{x}_N)]^\top$가 다변량 가우시안 분포를 따른다는 의미이다:
\[
\mathbf{f} \sim \mathcal{N}\bigl(\mathbf{m},\; \mathbf{K}\bigr)
\]
여기서 $\mathbf{m}_i = m(\mathbf{x}_i)$이고 $K_{ij} = k(\mathbf{x}_i, \mathbf{x}_j)$이다.

\textbf{GP 회귀}에서는 관측값 $y_i = f(\mathbf{x}_i) + \epsilon_i$, $\epsilon_i \sim \mathcal{N}(0, \sigma_n^2)$를 가정한다.
가우시안 우도 덕분에 사후 분포가 해석적으로 계산된다.
학습 데이터 $\mathcal{D} = \{(\mathbf{x}_i, y_i)\}_{i=1}^N$이 주어졌을 때,
새로운 입력 $\mathbf{x}_*$에서의 예측 분포는 다음과 같다:
\[
f_* \mid \mathcal{D}, \mathbf{x}_* \sim \mathcal{N}\bigl(\bar{f}_*,\; \mathrm{var}(f_*)\bigr)
\]
\begin{equation}
\bar{f}_* = \mathbf{k}_*^\top (\mathbf{K} + \sigma_n^2 \mathbf{I})^{-1} \mathbf{y}
\label{eq:gp_mean}
\end{equation}
\begin{equation}
\mathrm{var}(f_*) = k(\mathbf{x}_*, \mathbf{x}_*) - \mathbf{k}_*^\top (\mathbf{K} + \sigma_n^2 \mathbf{I})^{-1} \mathbf{k}_*
\label{eq:gp_var}
\end{equation}
여기서 $\mathbf{k}_* = [k(\mathbf{x}_1, \mathbf{x}_*), \dots, k(\mathbf{x}_N, \mathbf{x}_*)]^\top$이다.
GP 회귀의 핵심 장점은 예측값뿐 아니라 예측 불확실성을 동시에 제공한다는 점이다.


\subsection{GP 회귀에서 GP 분류로의 확장}

분류 문제에서 관측값 $y_i$는 이산적인 클래스 레이블이다.
GP 회귀의 가우시안 우도 $p(y \mid f) = \mathcal{N}(y \mid f, \sigma_n^2)$를 그대로 사용할 수 없다.
대신 잠재 함수(latent function) $f(\mathbf{x})$를 비선형 링크 함수로 변환하여 클래스 확률을 정의한다~\cite{Williams1998}.

\textbf{이진 분류}에서는 시그모이드(또는 probit) 링크를 사용한다:
\begin{equation}
p(y = 1 \mid f) = \sigma(f) = \frac{1}{1 + \exp(-f)}
\label{eq:sigmoid}
\end{equation}

\textbf{다중 클래스 분류} ($C \geq 3$)에서는 각 클래스 $c = 0, 1, \dots, C-1$에 대응하는
잠재 함수 $f_c(\mathbf{x})$를 도입하고, softmax 링크를 적용한다:
\begin{equation}
p(y = c \mid \mathbf{f}) = \frac{\exp(f_c)}{\sum_{c'=0}^{C-1} \exp(f_{c'})}
\label{eq:softmax}
\end{equation}
여기서 $\mathbf{f} = [f_0, f_1, \dots, f_{C-1}]^\top$이다.

이 비가우시안 우도로 인해 사후 분포 $p(\mathbf{f} \mid X, \mathbf{y})$는 해석적으로 계산되지 않는다.
이를 해결하기 위한 대표적인 근사 방법이 Laplace 근사와 Expectation Propagation (EP)이다~\cite{Nickisch2008}.
scikit-learn의 \texttt{GaussianProcessClassifier}는 Laplace 근사를 사용한다.


\subsection{Laplace 근사}

Laplace 근사는 사후 분포를 MAP(Maximum A Posteriori) 추정치 주변에서
2차 Taylor 전개하여 가우시안으로 근사하는 방법이다~\cite{MacKay1992}.

\textbf{사후 분포의 로그:}
\begin{equation}
\log p(\mathbf{f} \mid X, \mathbf{y}) = \log p(\mathbf{y} \mid \mathbf{f}) + \log p(\mathbf{f} \mid X) - \log p(\mathbf{y} \mid X)
\label{eq:log_posterior}
\end{equation}
여기서 $p(\mathbf{f} \mid X) = \mathcal{N}(\mathbf{f} \mid \mathbf{0}, \mathbf{K})$는 GP 사전 분포이다.

\textbf{Step 1: MAP 추정.}
비정규화 사후 분포를 최대화하는 $\hat{\mathbf{f}}$를 구한다:
\[
\hat{\mathbf{f}} = \arg\max_{\mathbf{f}} \; \Psi(\mathbf{f})
\]
여기서:
\begin{equation}
\Psi(\mathbf{f}) = \log p(\mathbf{y} \mid \mathbf{f}) - \frac{1}{2} \mathbf{f}^\top \mathbf{K}^{-1} \mathbf{f} - \frac{1}{2} \log |\mathbf{K}| - \frac{N}{2} \log(2\pi)
\label{eq:psi}
\end{equation}

이 최적화는 Newton--Raphson 반복으로 수행한다. $\Psi(\mathbf{f})$의 그래디언트는 다음과 같다:
\[
\nabla \Psi = \nabla \log p(\mathbf{y} \mid \mathbf{f}) - \mathbf{K}^{-1} \mathbf{f}
\]

\textbf{Step 2: Hessian 계산.}
MAP 추정치 $\hat{\mathbf{f}}$에서의 음의 Hessian을 계산한다:
\begin{equation}
\mathbf{W} = -\nabla\nabla \log p(\mathbf{y} \mid \mathbf{f})\big|_{\hat{\mathbf{f}}}
\label{eq:hessian}
\end{equation}

$\mathbf{W}$는 우도 함수의 곡률(curvature)을 나타내는 대각 행렬이다.
Softmax 우도의 경우, $\mathbf{W}$는 블록 대각 구조를 갖는다.

\textbf{Step 3: 가우시안 근사.}
사후 분포를 다음과 같이 근사한다:
\begin{equation}
q(\mathbf{f} \mid X, \mathbf{y}) = \mathcal{N}\bigl(\hat{\mathbf{f}},\; (\mathbf{K}^{-1} + \mathbf{W})^{-1}\bigr)
\label{eq:laplace_approx}
\end{equation}

근사 공분산 $(\mathbf{K}^{-1} + \mathbf{W})^{-1}$는 사전 분포의 공분산 $\mathbf{K}$보다 작다.
이는 관측 데이터가 잠재 함수에 대한 불확실성을 줄여주기 때문이다.


\subsection{새로운 입력에서의 예측}

학습 후, 새로운 입력 $\mathbf{x}_*$에서의 예측은 두 단계로 이루어진다.

\textbf{잠재 함수 예측.}
GP 사후를 이용하여 잠재 함수의 예측 분포를 구한다:
\[
f_* \mid \mathcal{D}, \mathbf{x}_* \sim q(f_*) = \mathcal{N}(\mu_*, \sigma_*^2)
\]
\begin{equation}
\mu_* = \mathbf{k}_*^\top \mathbf{K}^{-1} \hat{\mathbf{f}}
\label{eq:pred_mean}
\end{equation}
\begin{equation}
\sigma_*^2 = k(\mathbf{x}_*, \mathbf{x}_*) - \mathbf{k}_*^\top (\mathbf{K} + \mathbf{W}^{-1})^{-1} \mathbf{k}_*
\label{eq:pred_var}
\end{equation}

\textbf{클래스 확률 예측.}
잠재 함수 분포를 링크 함수에 적분하여 클래스 확률을 구한다:
\begin{equation}
p(y_* = c \mid \mathcal{D}, \mathbf{x}_*) = \int p(y_* = c \mid \mathbf{f}_*) \, q(\mathbf{f}_*) \, d\mathbf{f}_*
\label{eq:pred_class}
\end{equation}

이 적분은 해석적으로 풀리지 않으므로, 수치 적분 또는 Monte Carlo 근사를 사용한다~\cite{Neal1998}.
scikit-learn에서는 \texttt{max\_iter\_predict} 파라미터로 이 수치 적분의 반복 횟수를 제어한다.


\subsection{주변 우도와 하이퍼파라미터 최적화}

커널 함수의 하이퍼파라미터 $\boldsymbol{\theta}$를 학습 데이터로부터 결정하기 위해,
주변 우도(marginal likelihood)를 최대화한다:
\begin{equation}
\hat{\boldsymbol{\theta}} = \arg\max_{\boldsymbol{\theta}} \; \log p(\mathbf{y} \mid X, \boldsymbol{\theta})
\label{eq:hyperopt}
\end{equation}

Laplace 근사 하에서 주변 우도의 로그는 다음과 같이 근사된다:
\begin{equation}
\log p(\mathbf{y} \mid X, \boldsymbol{\theta}) \approx -\frac{1}{2} \hat{\mathbf{f}}^\top \mathbf{K}^{-1} \hat{\mathbf{f}} + \log p(\mathbf{y} \mid \hat{\mathbf{f}}) - \frac{1}{2} \log |\mathbf{K} + \mathbf{W}^{-1}|
\label{eq:marginal_likelihood}
\end{equation}

우변의 세 항은 각각 다음과 같이 해석된다.

\begin{table}[h]
\centering
\begin{tabular}{ll}
\toprule
항 & 의미 \\
\midrule
$-\frac{1}{2}\hat{\mathbf{f}}^\top \mathbf{K}^{-1}\hat{\mathbf{f}}$ &
  복잡도 벌칙 (complexity penalty). MAP 추정치가 사전 분포에서 \\
  & 벗어난 정도를 측정한다. \\
$\log p(\mathbf{y} \mid \hat{\mathbf{f}})$ &
  데이터 적합도 (data fit). 관측 레이블을 잘 설명하는 정도를 \\
  & 측정한다. \\
$-\frac{1}{2}\log|\mathbf{K} + \mathbf{W}^{-1}|$ &
  Occam 인자 (Occam factor). 모델 복잡도에 대한 자동 벌칙이다. \\
\bottomrule
\end{tabular}
\end{table}

이 세 항의 균형이 GP의 자동 복잡도 조절 능력을 제공한다.
과적합 모델은 데이터 적합도는 높지만 복잡도 벌칙이 크고, 과소적합 모델은 그 반대이다.

주변 우도의 $\boldsymbol{\theta}$에 대한 그래디언트를 계산하여 경사 기반 최적화를 수행한다.
다중 시작점(multi-start) 전략으로 지역 최적해에 빠지는 것을 완화한다.
본 구현에서는 \texttt{n\_restarts\_optimizer=3}으로 설정하여,
초기 하이퍼파라미터를 4번(1회 기본 + 3회 재시작) 최적화하고 가장 높은 주변 우도를 갖는 해를 선택한다.


\subsection{ARD RBF 커널}

본 구현에서 사용하는 커널은 상수 커널과 ARD(Automatic Relevance Determination) RBF 커널의 곱이다:
\begin{equation}
k(\mathbf{x}, \mathbf{x}') = \sigma_f^2 \exp\left(-\frac{1}{2} \sum_{d=1}^{D} \frac{(x_d - x'_d)^2}{\ell_d^2}\right)
\label{eq:ard_rbf}
\end{equation}
여기서:
\begin{itemize}[nosep]
  \item $\sigma_f^2$: 신호 분산(signal variance). 함수값의 전체적인 진폭을 제어한다.
  \item $\ell_d$: 차원 $d$의 길이 척도(length scale). 해당 차원에서의 함수 변동 규모를 결정한다.
\end{itemize}

\textbf{ARD의 물리적 해석.}
표준 RBF 커널은 모든 차원에 동일한 길이 척도 $\ell$을 사용한다.
ARD RBF 커널은 각 차원에 개별 길이 척도 $\ell_d$를 부여한다.
학습 결과 $\ell_d$가 크면, 차원 $d$의 변화가 함수값에 미치는 영향이 작다는 의미이다.
반대로 $\ell_d$가 작으면, 해당 차원의 변화에 함수가 민감하게 반응한다.

본 프로젝트의 3차원 파라미터 공간에서 ARD 길이 척도는 다음과 같이 해석된다:

\begin{table}[h]
\centering
\begin{tabular}{cccl}
\toprule
차원 $d$ & 파라미터 & 길이 척도 & 해석 \\
\midrule
0 & $\tilde{T}$ (정규화 추력) & $\ell_0$ & 추력 크기가 프로파일 유형에 미치는 영향 \\
1 & $\Delta a$ (반장축 변화) & $\ell_1$ & 궤도 크기 변화가 프로파일 유형에 미치는 영향 \\
2 & $\Delta i$ (경사 변화) & $\ell_2$ & 궤도면 변경이 프로파일 유형에 미치는 영향 \\
\bottomrule
\end{tabular}
\end{table}

학습된 길이 척도의 비율 $\ell_0 : \ell_1 : \ell_2$는 세 파라미터의 상대적 영향력(relevance)을 정량적으로 나타낸다.
이는 어떤 물리 파라미터가 추력 프로파일 유형 전환에 가장 큰 영향을 미치는지 식별하는 데 활용된다~\cite{Lee2023}.

\textbf{커널 파라미터의 초기값.}
코드에서 \texttt{ConstantKernel(1.0) * RBF(length\_scale=np.ones(n\_dims))}로 초기화한다.
상수 커널의 초기값은 $\sigma_f^2 = 1.0$이고, 모든 차원의 길이 척도 초기값은 $\ell_d = 1.0$이다.
정규화 좌표 $[0,1]^3$에서 길이 척도 1.0은 전체 범위에 해당하므로,
사전 정보가 없는 상태에서의 합리적인 출발점이다.


\subsection{One-vs-Rest (OvR) 다중 클래스 전략}

$C$개 클래스에 대한 다중 클래스 분류는 두 가지 대표적 전략으로 분해할 수 있다.

\textbf{One-vs-Rest (OvR).}
각 클래스 $c$에 대해 ``클래스 $c$ vs.\ 나머지 전체''로 이진 분류기를 독립적으로 학습한다.
총 $C$개의 이진 GP 분류기를 학습하며, 예측 시 각 분류기의 출력 확률을 정규화하여 클래스 확률을 구한다:
\begin{equation}
p(y = c \mid \mathbf{x}) = \frac{p_c(\mathbf{x})}{\sum_{c'=0}^{C-1} p_{c'}(\mathbf{x})}
\label{eq:ovr}
\end{equation}
여기서 $p_c(\mathbf{x})$는 클래스 $c$에 대한 이진 분류기의 양성 확률이다.

\textbf{One-vs-One (OvO).}
모든 클래스 쌍 $(c, c')$에 대해 이진 분류기를 학습한다. $\binom{C}{2}$개의 분류기가 필요하다.

본 구현에서는 scikit-learn의 \texttt{multi\_class='one\_vs\_rest'} 설정을 통해 OvR 전략을 사용한다.
$C = 3$일 때 OvR은 3개의 이진 GP 분류기를 학습한다.
이는 OvO의 3개 분류기와 같은 수이지만, OvR은 각 분류기가 전체 데이터를 사용하므로 데이터 효율성이 높다.

OvR 전략에서 각 이진 분류기는 독립적으로 Laplace 근사를 수행한다.
따라서 각 분류기는 자체적인 MAP 추정치 $\hat{\mathbf{f}}_c$와 Hessian $\mathbf{W}_c$를 가진다.
커널 하이퍼파라미터도 분류기별로 독립적으로 최적화된다.

그러나 OvR 전략에서 각 이진 분류기가 독립적으로 학습되므로,
원리적으로는 하이퍼파라미터(길이 척도)가 분류기마다 다를 수 있다.
scikit-learn 구현에서 \texttt{get\_length\_scales()} 메서드는 첫 번째(또는 합성된) 커널의 길이 척도를 반환한다.


\section{구현 매핑}

\noindent 대응 코드: \texttt{src/orbit\_transfer/sampling/gp\_classifier.py}

\subsection{클래스--수학 대응표}

\begin{table}[h]
\centering
\small
\begin{tabular}{llll}
\toprule
수학 개념 & 코드 요소 & 파일 위치 & 비고 \\
\midrule
GP 사전: $f \sim \mathcal{GP}(0, k)$ &
  \texttt{GaussianProcessClassifier(kernel=...)} &
  \texttt{L19--24} & zero-mean GP \\
ARD RBF 커널: $\sigma_f^2 \cdot \exp(\cdots)$ &
  \texttt{ConstantKernel * RBF(...)} &
  \texttt{L18} & 차원별 길이 척도 \\
Laplace 근사 &
  scikit-learn 내부 (기본값) &
  \texttt{L19} & 기본 추론 \\
OvR 다중 클래스 &
  \texttt{multi\_class='one\_vs\_rest'} &
  \texttt{L20} & 3개 이진 분류기 \\
예측 적분 반복 &
  \texttt{max\_iter\_predict=100} &
  \texttt{L21} & 수치 적분 반복 수 \\
다중 시작 최적화 &
  \texttt{n\_restarts\_optimizer=3} &
  \texttt{L22} & 총 4회 최적화 \\
학습: $\hat{\boldsymbol{\theta}} = \arg\max \log p(\mathbf{y} \mid X, \boldsymbol{\theta})$ &
  \texttt{fit(X, y)} &
  \texttt{L27--30} & 주변 우도 최대화 \\
확률 예측: $p(y_* = c \mid \mathcal{D}, \mathbf{x}_*)$ &
  \texttt{predict\_proba(X)} &
  \texttt{L32--34} & $(N, C)$ 배열 반환 \\
클래스 예측: $\hat{y}_* = \arg\max_c p(y_* = c)$ &
  \texttt{predict(X)} &
  \texttt{L36--38} & $(N,)$ 배열 반환 \\
ARD 길이 척도: $[\ell_0, \ell_1, \ell_2]$ &
  \texttt{get\_length\_scales()} &
  \texttt{L40--51} & 커널 파라미터 추출 \\
\bottomrule
\end{tabular}
\end{table}


\subsection{\texttt{GPClassifierWrapper} 클래스 구조}

클래스는 다음 네 가지 공개 메서드를 제공한다.

\textbf{\texttt{\_\_init\_\_(self, n\_dims=3)}} (lines 16--25).
커널과 분류기를 초기화한다. \texttt{n\_dims}는 입력 차원수로, 기본값 3은 파라미터 공간의 차원에 대응한다.
커널을 \texttt{ConstantKernel(1.0) * RBF(length\_scale=np.ones(n\_dims))}로 구성하여 ARD RBF 커널을 정의한다.
scikit-learn의 \texttt{Product} 커널 객체가 생성되며, \texttt{k1}이 \texttt{ConstantKernel}, \texttt{k2}가 \texttt{RBF}에 해당한다.

\textbf{\texttt{fit(self, X, y)}} (lines 27--30).
학습 데이터 $X \in \mathbb{R}^{N \times d}$와 레이블 $y \in \{0, 1, 2\}^N$으로 GP 분류기를 학습한다.
내부적으로 scikit-learn이 다음을 수행한다:
\begin{enumerate}[nosep]
  \item OvR 전략에 따라 3개의 이진 GP 분류기를 구성
  \item 각 분류기에 대해 Laplace 근사로 MAP 추정
  \item 주변 우도를 최대화하여 커널 하이퍼파라미터 최적화
  \item \texttt{n\_restarts\_optimizer=3}에 따라 4회 반복, 최적 결과 선택
\end{enumerate}

학습 후 \texttt{self.is\_fitted = True}로 설정하고, 학습된 커널은 \texttt{self.gpc.kernel\_} 속성에 저장된다.
초기 커널 \texttt{self.gpc.kernel}과 구분된다는 점에 유의해야 한다.
밑줄(\texttt{\_})이 붙은 속성이 학습 결과이다.

\textbf{\texttt{predict\_proba(self, X)}} (lines 32--34).
새로운 입력 $X \in \mathbb{R}^{N_* \times d}$에서 클래스 확률 행렬 $(N_*, C)$를 반환한다.
각 행은 확률 벡터로, 합이 1이다.
이 확률 분포는 능동 학습의 predictive entropy 계산에 직접 사용된다.

\textbf{\texttt{get\_length\_scales(self)}} (lines 40--51).
학습된 ARD 길이 척도 $[\ell_0, \ell_1, \ell_2]$를 반환한다. 학습 전이면 \texttt{None}을 반환한다.
scikit-learn의 \texttt{Product} 커널에서 \texttt{kernel\_.k2}가 RBF 부분이므로,
이로부터 \texttt{length\_scale} 속성을 추출한다.
\texttt{hasattr} 검사로 커널 구조가 예상과 다른 경우에도 안전하게 동작한다.


\subsection{scikit-learn 내부 구현 요약}

\texttt{GPClassifierWrapper}는 scikit-learn~\cite{Pedregosa2011}의 \texttt{GaussianProcessClassifier}를 직접 래핑한다.
scikit-learn의 내부 동작을 요약하면 다음과 같다.

\textbf{학습 과정 (\texttt{fit}):}
\begin{enumerate}[nosep]
  \item \texttt{multi\_class='one\_vs\_rest'}이면, \texttt{OneVsRestClassifier}가 클래스 수 $C$개의 이진 \texttt{\_BinaryGaussianProcessClassifierLaplace} 객체를 생성한다.
  \item 각 이진 분류기에서:
  \begin{itemize}[nosep]
    \item 커널 행렬 $\mathbf{K}$를 계산한다 ($O(N^2 d)$).
    \item Newton--Raphson으로 MAP 추정치 $\hat{\mathbf{f}}$를 구한다.
    \item Hessian $\mathbf{W}$를 계산한다.
    \item 주변 우도의 그래디언트를 이용하여 L-BFGS-B로 하이퍼파라미터를 최적화한다.
    \item \texttt{n\_restarts\_optimizer}만큼 랜덤 초기점에서 반복하여 최적 해를 선택한다.
  \end{itemize}
\end{enumerate}

\textbf{예측 과정 (\texttt{predict\_proba}):}
\begin{enumerate}[nosep]
  \item 각 이진 분류기에서 잠재 함수의 예측 평균 $\mu_*$과 분산 $\sigma_*^2$를 계산한다.
  \item 시그모이드 적분을 수치적으로 수행하여 양성 확률 $p_c(\mathbf{x}_*)$를 구한다.
  \item OvR 정규화: $p(y = c) = p_c / \sum_{c'} p_{c'}$를 적용한다.
\end{enumerate}

\textbf{계산 복잡도.}
학습 시 커널 행렬의 역행렬 계산이 지배적이다.
시간 복잡도는 $O(N^3)$, 공간 복잡도는 $O(N^2)$이다.
$N = 500$ (본 프로젝트의 최대 샘플 수) 수준에서는 수 초 이내에 학습이 완료된다.


\section{수치 검증}

\noindent 대응 테스트: \texttt{tests/test\_sampling.py::TestGPClassifier}

\subsection{합성 3-class 분류 정확도}

\begin{lstlisting}
class TestGPClassifier:
    def test_synthetic_3class(self):
        """Synthetic 3-class data: accuracy > 80%."""
        rng = np.random.default_rng(42)
        N = 200
        X = rng.uniform(0, 1, (N, 3))
        # Simple rule: class = 0 if x0<0.33, 1 if 0.33<x0<0.66, 2 otherwise
        y = np.where(X[:, 0] < 0.33, 0,
                     np.where(X[:, 0] < 0.66, 1, 2))

        gpc = GPClassifierWrapper(n_dims=3)
        # 80% train, 20% test
        n_train = 160
        gpc.fit(X[:n_train], y[:n_train])
        pred = gpc.predict(X[n_train:])
        accuracy = np.mean(pred == y[n_train:])
        assert accuracy > 0.80, f"accuracy {accuracy:.2f} < 0.80"
\end{lstlisting}

이 테스트는 다음 사항을 검증한다.

\textbf{데이터 생성.}
$[0,1]^3$ 공간에서 200개 점을 균일 분포로 추출한다.
클래스 레이블은 첫 번째 차원 $x_0$에 의해서만 결정된다:
$x_0 < 0.33$이면 class 0, $0.33 \leq x_0 < 0.66$이면 class 1, $x_0 \geq 0.66$이면 class 2이다.
이는 결정 경계가 $x_0 = 0.33$과 $x_0 = 0.66$인 단순한 구조이다.

\textbf{학습/테스트 분할.}
160개로 학습, 40개로 테스트한다 (80/20 분할).

\textbf{정확도 기준.}
테스트 정확도 80\% 이상을 요구한다.
ARD RBF 커널은 학습 과정에서 $\ell_0$가 작아지고 $\ell_1, \ell_2$가 커지는 방향으로 수렴하여,
$x_0$ 차원의 지배적 영향을 학습해야 한다.

\textbf{검증 의미.}
이 테스트는 (1) \texttt{fit}--\texttt{predict} 파이프라인이 정상 동작하고,
(2) ARD 커널이 관련 차원을 식별할 수 있으며,
(3) Laplace 근사가 충분한 분류 성능을 제공함을 확인한다.


\subsection{길이 척도 추출 검증}

합성 테스트에서 학습된 ARD 길이 척도를 \texttt{get\_length\_scales()}로 추출할 수 있다.
위의 데이터 구조에서, 클래스가 $x_0$에 의해서만 결정되므로, 학습 후 다음 관계가 기대된다:
\[
\ell_0 \ll \ell_1 \approx \ell_2
\]

$\ell_0$가 작다는 것은 함수가 $x_0$ 방향으로 빠르게 변한다는 의미이며,
$\ell_1, \ell_2$가 크다는 것은 $x_1, x_2$ 방향의 변화가 함수에 거의 영향을 미치지 않는다는 의미이다.
이 관계는 ARD의 자동 관련성 결정 능력을 보여준다.


\subsection{미학습 상태 안전성}

\texttt{get\_length\_scales()} 메서드는 학습 전 (\texttt{is\_fitted=False}) 호출 시 \texttt{None}을 반환한다.
이는 학습되지 않은 커널 파라미터에 접근하는 오류를 방지한다.


\subsection{확률 출력 일관성}

\texttt{predict\_proba(X)}의 반환값은 $(N, C)$ 형상의 배열이며, 각 행의 합은 1이다.
OvR 전략에서 각 이진 분류기의 양성 확률을 정규화하므로, 이 성질이 보장된다:
\[
\sum_{c=0}^{C-1} p(y = c \mid \mathbf{x}) = 1, \quad \forall \mathbf{x}
\]


\subsection{능동 학습 맥락에서의 통합 검증}

\texttt{tests/test\_sampling.py::TestAdaptiveSampler::test\_small\_convergence} 테스트는
\texttt{GPClassifierWrapper}가 능동 학습 루프 내에서 반복적으로 학습--예측되는 전체 파이프라인을 검증한다.
20개의 초기 LHS 샘플에서 시작하여 최대 50개까지 적응적으로 샘플을 추가하며,
GP 분류기가 매 반복마다 재학습된다.
이 테스트는 GP 분류기가 점진적 데이터 추가에 대해 안정적으로 동작하고,
predictive entropy 기반 acquisition function과 올바르게 연동됨을 확인한다.


\bibliographystyle{IEEEtran}
\bibliography{references}

\end{document}
